\chapter{CD4 Lymphocyte Count}

\section*{What Is This Test?}
A CD4 lymphocyte count measures the absolute number of CD4+ T-lymphocytes per microliter of blood. CD4+ T cells are the primary target and reservoir of HIV, and their progressive depletion drives the immunodeficiency that defines AIDS. The CD4 count is used to stage HIV disease, decide when to start antiretroviral therapy (ART), initiate prophylaxis against opportunistic infections (PCP, MAC), and monitor immunologic response to ART. CD4 is also used in other immunodeficiency evaluations (chronic corticosteroid therapy, hematologic malignancy, transplant recipients).

\section*{Alternative Names}
CD4 Count; CD4+ T-Lymphocyte Count; T-Helper Cell Count; CD4 Absolute Count; T-Cell Subset Panel

\section*{Principle}
Whole blood collected in EDTA is stained with fluorescent monoclonal antibodies against CD3, CD4, and CD8 surface markers and analyzed by flow cytometry. A known number of fluorescent beads or the volumetric count provides an absolute CD4 count per microliter. Typical panels report absolute CD4, absolute CD8, CD4/CD8 ratio, and percent lymphocyte subsets. Single-platform flow cytometry uses absolute count beads; dual-platform uses total WBC and differential from the hematology analyzer multiplied by CD4\%.

\section*{History}
CD4 became a clinical marker of AIDS soon after the recognition of the syndrome in 1981 and the identification of HIV as the cause in 1983. CD4 thresholds for PCP prophylaxis ($<$ 200) were established by the late 1980s; initial ART guidelines used CD4 $<$ 500, evolving to universal ART regardless of CD4 in 2015--2016.

\section*{Why This Test Is Ordered}
\begin{itemize}
  \item Initial staging of HIV infection after a positive HIV confirmatory test
  \item Decision to initiate opportunistic infection prophylaxis (CD4 $<$ 200: PCP; $<$ 50: MAC)
  \item Monitor immunologic response after starting ART (expect rise of 50--150 cells/$\mu$L in first year)
  \item Risk assessment in chronic immunosuppression (transplant, corticosteroids, chemotherapy)
  \item Evaluation of unexplained immunodeficiency, lymphopenia, or opportunistic infection
\end{itemize}

\section*{Is This a Primary or Secondary Test}
Secondary test; follows positive HIV diagnosis for staging and monitoring.

\section*{How This Test Is Performed}
Blood drawn in EDTA (lavender-top tube) is transported to the flow cytometry laboratory and processed within 24--72 hours. Timing matters: specimens beyond 72 hours may have falsely low counts. Smoking or recent exercise minimally affect results. Diurnal variation is small (lowest in morning).

\section*{What Preparation Is Needed}
None. Measure before blood transfusion if possible.

\section*{Contraindications and Precautions}
Specimens must be at room temperature and analyzed within 72 hours. Timing relative to vaccines and acute illness may transiently affect the count. Recent corticosteroid therapy lowers CD4 counts.

\section*{Risks}
None beyond venipuncture.

\section*{What Happens After the Test}
Results are reported as absolute CD4 count per microliter, percent CD4 of lymphocytes, and CD4/CD8 ratio. Initially every 3--4 months in untreated patients and confirms of therapy; at 3 and 6 months after starting or changing ART, then every 6--12 months once stable on ART and undetectable viral load.

\section*{Understanding Results}

\subsection*{Below Normal}
\begin{itemize}
  \item Mild depletion: 200--499 cells/$\mu$L -- increased risk of candidiasis, Kaposi sarcoma
  \item Severe depletion: $<$ 200 cells/$\mu$L -- AIDS-defining; start PCP prophylaxis (TMP-SMX)
  \item Very severe: $<$ 50 cells/$\mu$L -- start MAC prophylaxis (azithromycin); high risk of CMV, cryptococcal meningitis, disseminated MAC
  \item Other causes of low CD4: corticosteroid use, chemotherapy, radiation, marrow failure; severe sepsis; congenital immunodeficiencies (DiGeorge, IPEX)
\end{itemize}

\subsection*{Above Normal}
\begin{itemize}
  \item Values $>$ 1,500/$\mu$L are physiologic in infants and small children; expected recovery 200--1,500 in adults with virologic suppression
  \item Persistently high CD4 is rarely pathological; consider lymphoproliferative disorders if concurrent lymphocytosis
\end{itemize}

\section*{Normal Results}
\begin{itemize}
  \item \textbf{Adult CD4 absolute:} 500--1,500 cells/$\mu$L
  \item \textbf{CD4percent:} 30--60\% of lymphocytes
  \item \textbf{CD4/CD8 ratio:} 1.0--4.0 (typical $>$ 1.0)
  \item \textbf{Infants/children:} Absolute CD4 substantially higher, percent $>$ 25\%
  \item \textbf{AIDS-defining threshold:} $<$ 200 cells/$\mu$L or CD4\% $<$ 14\%
\end{itemize}

\section*{Follow-up Tests}
\begin{itemize}
  \item \textbf{HIV viral load (HIV-1 RNA):} Monitors virologic suppression; expect $<$ 50 copies/mL by 24 weeks on ART
  \item \textbf{HIV genotype resistance test:} At baseline and treatment failure
  \item \textbf{CBC with differential:} Concurrent lymphopenia or anemia
  \item \textbf{Pneumocystis prophylaxis:} Begin when CD4 $<$ 200 cells/$\mu$L
  \item \textbf{MAC prophylaxis:} Begin azithromycin when CD4 $<$ 50 cells/$\mu$L
  \item \textbf{Quantiferon-TB Gold and syphilis, hepatitis B/C serologies:} At baseline HIV evaluation
\end{itemize}

\section*{References}
\begin{enumerate}
  \item Panel on Antiretroviral Guidelines for Adults and Adolescents. Guidelines for the Use of Antiretroviral Agents in Adults and Adolescents with HIV. U.S. Department of Health and Human Services; updated 2023.
  \item Maldarelli F, Palmer S. HIV. In: Bennett JE, et al., eds. Mandell, Douglas, and Bennett's Principles and Practice of Infectious Diseases. 9th ed. Elsevier; 2020.
  \item Clinical and Laboratory Standards Institute. Enumeration of Immunologically Defined Cell Populations by Flow Cytometry. CLSI Guideline H42-A3; 2022.
  \item Centers for Disease Control and Prevention. Revised surveillance case definition for HIV infection --- United States, 2014. MMWR. 2014;63(RR-03):1-10.
\end{enumerate}

\clearpage